\begin{resumo}[Abstract]
 \begin{otherlanguage*}{english}
   This work investigates the low test coverage in AgroMart Web and its impacts on maintainability and on the safe evolution of the application.
   The initial diagnosis, supported by SonarQube, identified coverage of only 0.5\%, with 2,186 eligible lines and 14 lines covered by tests,
   indicating a high risk of regressions during code changes. Given this scenario, the objective of the study is to improve the project's internal
   quality through the progressive expansion of unit tests in fundamental system units. The adopted methodology is exploratory, structured in cycles of diagnosis,
   unit prioritization, test implementation, metric assessment, and consolidation of practices in development documentation.
   As a result, the work establishes a technical baseline and an evidence-driven continuous improvement flow, integrating static analysis,
   test execution, and systematic monitoring of quality indicators. It is concluded that the proposed strategy provides a feasible path
   to sustainably increase test coverage and reduce maintenance risks in the context of AgroMart Web's collaborative evolution.

   \vspace{\onelineskip}
 
   \noindent 
   \textbf{Key-words}: software quality. test coverage. unit testing. AgroMart.
 \end{otherlanguage*}
\end{resumo}
